%%%%%%%%%%%%%%%%%%%%%%%%%%%%%%%%%%%%%%%%%%%%%%%%%%%%%%%%%%%
% Speaker notes for the PPI presentation.
% Standalone reference document, prints as a compact sheet.
% Slide order matches slides.tex.
%%%%%%%%%%%%%%%%%%%%%%%%%%%%%%%%%%%%%%%%%%%%%%%%%%%%%%%%%%%

\documentclass[11pt,a4paper]{article}
\usepackage[margin=2cm]{geometry}
\usepackage{enumitem}
\usepackage{xcolor}

\setlist[itemize]{itemsep=2pt,topsep=2pt,leftmargin=*}
\setlength{\parindent}{0pt}
\setlength{\parskip}{0.4em}

\definecolor{myaccent}{HTML}{C0306E}
\newcommand{\slidenote}[1]{\par\vspace{0.8em}{\large\bfseries\color{myaccent}#1}\par\vspace{0.2em}}

\begin{document}

{\Large\bfseries Speaker notes \hfill PPI Endometriosis, MIMUW 2026}\par
\rule{\textwidth}{0.4pt}

\slidenote{Data and network construction}
\begin{itemize}
  \item DisGeNET: curated gene-disease database
  \item C0014175 = UMLS endometriosis ID
  \item STRING: physical and functional protein interactions
  \item First-shell expansion: adds direct neighbours, cap 300
  \item Density driven by the expansion
\end{itemize}

\slidenote{Null models}
\begin{itemize}
  \item ER, Erd\H{o}s-R\'enyi: random edges at endo-ppi's density; tests ``edge count alone?''
  \item BA, Barab\'asi-Albert: preferential attachment with $m \approx 53$, -> so total edge count matches PPI; tests ``scale-free distribution alone?''
  \item Pref attachment: rich-get-richer; each new node adds $m$ edges to existing nodes with probability $\propto$ degree, producing a scale-free distribution
  \item Scale-free: power-law degree distribution $P(k) \propto k^{-\alpha}$; many low-degree nodes plus a few high-degree hubs, no typical degree
  \item CM, configuration model: preserves degree sequence; tests ``degrees alone?''
  \item CM by double-edge swap: pick two edges, swap their endpoints; preserves every node's degree while randomizing topology
  \item \texttt{nx.configuration\_model} would lose 20\% of edges in dense graph, due to parallel edges and self-loops collapsing on simplification
\end{itemize}

\slidenote{Network statistics}
\begin{itemize}
  \item Clustering: fraction of a node's neighbor pairs that are also connected
  \item PPI 0.637 vs CM 0.512, same degree sequence; 25\% gap, so ppi not so random
  \item Assortativity: degree correlation at edge endpoints
  \item PPI uniquely positive, $+0.052$; hubs cluster together
  \item CM disassortative, $-0.181$; too few hubs to pair with each other randomly
\end{itemize}

\slidenote{Degree distribution}
\begin{itemize}
  \item PPI heavy-tailed: low-degree nodes + extreme hubs, highest = 100 * lowest  -> span two orders of magnitude
  \item ER: narrow Poisson peak around mean 93, no tail
  \item BA: power-law shape by design but scattered here, minimum degree fixed at $m \approx 53$
  \item CM matches PPI by construction
\end{itemize}

\slidenote{Diffusion models}
\begin{itemize}
  \item SIR: each step infected node infect each neighbor with prob $\beta$, then recovers with prob $\gamma$
  \item IC: newly activated node gets one chance per inactive neighbor with prob $p$, then exhausted
  \item Both are Monte Carlo simulations, 50 runs per config, results averaged
  \item Pairing the two checks whether results depend on choice of dynamics
\end{itemize}

\slidenote{Calibration sweep}
\begin{itemize}
  \item Too high: network saturates, every blocking strat looks ineffective
  \item Too low: propagation dies, every strat looks artificially good
  \item Target: 60--85\% of the network without blocking
  \item Dense graph forces pretty low parameters.
  \item Each row = mean over 50 Monte Carlo runs at fixed parameters
\end{itemize}

\slidenote{Group-centrality algorithms}
\begin{itemize}
  \item Structural: look only at graph topology, no diffusion simulation
  \item ``Lazy'' in CELF = recompute spread only for candidates that could be stale; cuts the cost by orders of magnitude
  \item the problem of finding truly optimal set is NP-hard
  \item gdeg attacks a submodular function, so greedy achieves at least $(1 - 1/e)$ of the optimum (Nemhauser et al. 1978)
  \item for each k total 4 result chains: gdeg, gbtw, CELF + SIR, CELF + IC
\end{itemize}

\slidenote{Top 10 picks per chain}
\begin{itemize}
  \item gdeg \& gbtw share IL6, ESR1, EGFR, INS in top-10: topological quality
  \item ESR1 is the estrogen receptor, biologically central to endometriosis
  \item \texttt{celf\_sir} top picks are growth-factor / angiogenesis proteins (FGF4, KDR, PDGFRB)
  \item \texttt{celf\_ic} top picks are immune-modulation proteins (CD40, LAG3, CX3CL1)
  \item CELF chains share 0 with each other and 0 with structural chains in top-10
  \item gdeg-gbtw overlap of 4/10 is well above chance; everything else at chance baseline
\end{itemize}

\slidenote{Spread reduction across algorithms}
\begin{itemize}
  \item gbtw plateau at $k \approx 200$, hits 98.7\% SIR / 99.1\% IC
  \item Next-best (celf\_ic) behind by 12--17 pp at $k = 200$
  \item Random baseline lowest at every $k \ge 100$
  \item CELF underperforms despite optimizing the spread metric directly
  \item Why CELF loses: dense PPI rewards high-betweenness hubs
\end{itemize}

\slidenote{PPI vs null networks}
\begin{itemize}
  \item PPI vs CM (same degree sequence) under gbtw: 7.8 pp gap = real higher-order structure
  \item BA vulnerable to gdeg because extreme hubs fragment the graph fast
  \item PPI doesn't show that vulnerability despite heavy tail; clustering buffers it
\end{itemize}

\slidenote{Top-20 composition}
\begin{itemize}
  \item Union of top-20 across 4 chains: 75 distinct proteins
  \item Only 5 appear in 2+ chains, all gdeg-gbtw: ESR1, IL6, EGFR, INS, CYP3A4
  \item Every other pairwise comparison: zero overlap in top-20
  \item 93\% algorithm-specific = answer is dominated by which algorithm you ask
  \item celf\_sir cluster: angiogenesis + ECM remodeling (FGF4, KDR, PDGFRB, VTN, TIMP1)
  \item celf\_ic cluster: immune modulation + transcription (CD40, LAG3, HDAC2, MED4)
\end{itemize}

\slidenote{Conclusion and limitations}
\begin{itemize}
  \item Best $k$-set: prefix of gbtw chain, saturates around $k = 200$ (44\% of nodes)
  \item Strongly algorithm-dependent: 93\% of top-20 algorithm-specific, CELF chains share 0 with structural
  \item At realistic small $k$ (10 proteins), even best algorithm achieves only ~7\% reduction
  \item Diffusion-aware chains surface angiogenesis (SIR) and immune-modulation (IC) targets that structural methods miss
\end{itemize}

\end{document}
